\chapter{Conclusions and Future Directions}
\label{ch:conclusion}

\section{Conclusions}
\label{sec:conclusions}

Modern video retrieval must operate across heterogeneous query forms, long and informationally dense videos, and candidate collections whose scale precludes exhaustive scoring. This thesis has argued that these demands are best understood as three sequential scaling challenges along the retrieval pipeline: semantic scaling, observational scaling, and search scaling. Each challenge corresponds to a distinct stage of the pipeline, and each presents a distinct bottleneck that is not eliminated by optimizing an upstream stage alone.

This thesis is organized around three studies, each addressing one of these stages. The first study establishes the representation foundation through semantic scaling. The second extends the investigation to observational scaling, studying how visual evidence should be allocated for long and multi-granular video retrieval. The third addresses search scaling, examining the retrieval infrastructure required when learned relevance signals are deployed at web scale. Together, they support a shared thesis: scalable video retrieval requires structure-aware scaling across semantic capability, visual observation, and search execution, and such scaling should be evaluated through generalization beyond the conditions under which each stage is optimized. The key findings and contributions of each study are summarized below.

\subsection{Summary of Contributions}

\paragraph{Chapter~\ref{ch:universal-video-retrieval}: Universal Video Retrieval through Capability-Oriented Learning.}
Existing video retrieval research often evaluates on a narrow set of benchmarks, leaving critical capability gaps unexamined. This work proposed the General Video Embedder (GVE) framework, which (1) identified deficiencies in fine-grained, temporal, and partially relevant retrieval through a structured assessment across 16 public datasets; (2) introduced V-SynFlow, a synthesis pipeline that produced the Universal Video Retrieval Dataset (UVRD) with over 1.55 million multi-task video-text pairs targeting those gaps; and (3) developed a Modality Pyramid curriculum that leverages inherent task hierarchies for progressive learning. The resulting GVE model achieved state-of-the-art zero-shot performance across diverse retrieval scenarios, demonstrating that conventional benchmarks poorly correlate with holistic retrieval capability and that partially relevant retrieval is a pivotal yet overlooked real-world challenge.

\paragraph{Chapter~\ref{ch:spatiotemporal-scaling}: Spatiotemporal Scaling Analysis and Efficient Inference.}
The way videos are observed during embedding generation fundamentally affects retrieval quality, yet this dimension has received limited attention. This work introduced the Multi-Granular Long Video Retrieval (MG-LVR) benchmark, the first benchmark supporting multi-granular queries for long-form video retrieval. Evaluation across 23 multimodal embedding models revealed a significant granularity gap---fine-grained retrieval lags coarse-grained performance by over 75\% under sparse sampling. The comprehensive spatiotemporal scaling analysis identified actionable principles including frame-count anchoring and granularity-dependent token allocation. Building on these insights, the Video Embedding Waterfall (VEterfall) framework was proposed, achieving up to 20\% improvement in fine-grained query performance while operating in 4$\times$ to 16$\times$ more efficient inference settings.

\paragraph{Chapter~\ref{ch:scale-efficient-retrieval}: Scale-Efficient Retrieval via Co-Designed Metrics, Graph Search, and GPU Execution.}
A learned retrieval signal must not only be accurate but also searchable and deployable at scale. This work proposed the GPU-accelerated Multi-relational Parallel Graph Retrieval (GMP-GR) framework with three key innovations: (1) Multi-relational User-Item Relevance Metric Learning (MUIRML) that unifies diverse user behaviors through multi-objective optimization and employs a self-covariant loss to enhance graph pathfinding; (2) Hierarchically Parallel Approximate Nearest Neighbor Search (HiPANNS) that conducts breadth-depth-balanced graph traversal with adaptive GPU aggregation for irregular neural computation; and (3) system-level optimizations including traffic dispersal, compilation optimization, model quantization, and timely stream update. Deployed across more than twenty applications at Baidu, GMP-GR serves hundreds of millions of users with throughput exceeding one hundred million requests per second, improving more than 3\% time of page for users.

\subsection{Unified Methodology and Insights}

Together, these studies form a coherent methodology that extends from representation learning through observation-aware embedding generation to deployment-scale retrieval. Each study addresses the scaling bottleneck at its respective stage of the pipeline, and each depends on the conclusions of the preceding one: the first study establishes what capabilities a retrieval model should possess, the second examines how visual observation affects the realization of those capabilities, and the third ensures that the resulting retrieval signals can scale to large candidate collections.

A central insight shared across all three studies is that scaling capacity does not automatically translate into retrieval generalization. Scaling model size alone does not cover the full breadth of semantic capabilities. Increasing observation density does not automatically improve fine-grained retrieval. Improving metric accuracy does not guarantee faster search. Within the scope of this thesis, generalizability serves as an empirical criterion for assessing whether scaling has been performed in a structure-aware manner at each stage of the retrieval pipeline: only by understanding the structure of the bottleneck and designing accordingly does the system remain effective under conditions it was not explicitly optimized for. This unified view provides practical guidance for building video retrieval systems that support richer semantic tasks, adapt to different sampling budgets, and scale to large candidate spaces while maintaining effectiveness across heterogeneous queries, videos, and deployment conditions.

The notion of generalizability evolves across the three stages of the retrieval pipeline. In semantic scaling, it concerns whether the learned representation transfers across query modalities, retrieval domains, and relevance granularities. In observational scaling, it concerns whether retrieval quality remains robust when visual evidence is resampled across different densities, durations, and query granularities. In search scaling, it concerns whether the learned retrieval signal remains effective under different deployment contexts, traffic patterns, and serving constraints. These forms are not identical, but they are unified by the same principle: scaling is meaningful only when the resulting capability transfers beyond the narrow condition under which it is optimized.

\section{Future Directions}
\label{sec:future-works}

The three studies in this thesis address distinct dimensions of video retrieval, yet each opens avenues that extend beyond the scope of the present work. This section outlines four directions for future research, organized from immediate extensions of the current findings to longer-term visions for the field.

\subsection{Query- and Content-Adaptive Observation}

The results in this thesis suggest that different queries do not require the same visual evidence. Coarse topical queries, fine-grained event queries, visual queries, and composed queries may prefer different allocations of temporal coverage and spatial detail. A natural next step is to make the observation policy adaptive to both the query and the video content, so that the system can spend computation where evidence is most likely to matter.

One promising direction is to develop lightweight query classifiers that predict the optimal sampling configuration before video encoding. Such classifiers could analyze query text to estimate temporal granularity, spatial specificity, and relevance distribution patterns, then select appropriate frame counts and resolutions accordingly. Another direction is content-aware sampling, where the system analyzes video content to identify information-dense segments that warrant higher sampling density. This could involve preliminary passes with cheap feature extractors to detect scene changes, action boundaries, or semantically rich regions, followed by targeted dense sampling in those regions.

More ambitiously, end-to-end differentiable sampling policies could be trained jointly with the retrieval model, allowing the system to learn task-specific observation strategies through reinforcement learning or attention mechanisms. Such approaches could enable dynamic allocation of computational resources within a single video, concentrating frames on moments most relevant to the query while maintaining sufficient coverage to avoid missing distributed evidence.

\subsection{Retrieval-Oriented Video Foundation Models}

Current multimodal models are often trained for broad vision-language instruction following or contrastive alignment. Future video foundation models could be trained more explicitly around retrieval requirements, including partial relevance, temporal evidence preservation, multimodal query support, robustness under sparse observation, and compatibility with efficient search.

The findings from Chapter~\ref{ch:spatiotemporal-scaling} suggest that current video models exhibit frame-count anchoring, where temporal representations become overly specialized to specific training densities. Future training objectives could explicitly encourage density-invariant representations, perhaps through multi-density contrastive learning where the same video must retrieve the same query across different sampling configurations. Similarly, the capability gaps identified in Chapter~\ref{ch:universal-video-retrieval} suggest that training data should more systematically cover fine-grained temporal reasoning, spatial relation understanding, and composed query scenarios.

Another important direction is developing video foundation models that natively support partial relevance, where a video is relevant to a query even if only a portion of its content matches. This requires moving beyond global video-level representations toward segment-aware or moment-aware embeddings that can capture localized relevance. Such models might employ hierarchical encoding schemes that maintain both coarse global summaries and fine-grained local representations, enabling retrieval at multiple granularity levels simultaneously.

Furthermore, future video foundation models should be designed with deployment constraints in mind. This includes compatibility with approximate search structures, efficient update mechanisms for dynamic corpora, and representations that remain discriminative even under aggressive compression or quantization for large-scale storage.

\subsection{Joint Optimization of Embeddings and Indexes}

Large-scale retrieval depends on both representation quality and index searchability. Future work can further close the loop between embedding learning and index behavior, optimizing models not only for pairwise similarity or ranking accuracy, but also for graph navigability, update efficiency, memory footprint, batching behavior, and hardware execution.

The work in Chapter~\ref{ch:scale-efficient-retrieval} demonstrated that metric properties like approximate triangle inequality significantly impact graph traversal efficiency. Future research could explore learned metrics that are explicitly optimized for specific index structures, such as hierarchical navigable small world graphs or product quantization schemes. This could involve differentiable approximations of graph traversal that allow gradient-based optimization of embedding spaces for both retrieval accuracy and search efficiency.

Another promising direction is co-designing embeddings and indexes for dynamic scenarios where the candidate corpus evolves over time. This requires representations that support efficient incremental indexing without full rebuilds, as well as metrics that remain stable as the corpus distribution shifts. Online learning approaches that continuously update both embeddings and index structures could enable more responsive retrieval systems for applications like news recommendation or trending content discovery.

Hardware-aware embedding learning represents another frontier. Different hardware platforms (CPUs, GPUs, TPUs, specialized accelerators) have distinct computational characteristics that affect retrieval performance. Future work could develop embedding models that are co-optimized with specific hardware constraints, such as memory bandwidth, parallelism patterns, or instruction set capabilities. This could lead to significant efficiency gains in large-scale deployment scenarios.

\subsection{Evaluation Protocols for Real-World Video Retrieval}

Existing benchmarks remain useful, but real retrieval systems involve changing corpora, mixed query modalities, variable video duration, latency budgets, freshness constraints, and downstream ranking interactions. Building evaluation protocols that reflect these conditions will be essential for measuring progress beyond static benchmark recall.

The benchmarks developed in this thesis, including UVRD and MG-LVR, provide foundations for more comprehensive evaluation, but significant gaps remain. Future evaluation frameworks should incorporate temporal dynamics, where videos are added, removed, or modified over time, testing the system's ability to maintain retrieval quality in evolving corpora. They should also include diverse query modalities within the same evaluation, requiring systems to handle text, image, video, and composed queries seamlessly.

Latency and throughput constraints should be explicitly incorporated into evaluation metrics, moving beyond pure recall measurements to efficiency-adjusted metrics that reward systems achieving high accuracy within computational budgets. This is particularly important for real-time applications where response time directly impacts user experience.

Downstream task evaluation represents another important direction. Rather than evaluating retrieval in isolation, future protocols should measure how retrieval quality affects downstream tasks like recommendation, content discovery, or user engagement. This requires end-to-end evaluation frameworks that integrate retrieval with ranking, filtering, and presentation components.

Finally, evaluation should consider fairness, diversity, and serendipity aspects of retrieval, ensuring that systems do not merely optimize for narrow relevance metrics but also support broader information access goals. This includes evaluating retrieval across different user populations, content categories, and temporal contexts to ensure robust and equitable performance.
